\documentclass[sigconf,anonymous]{acmart} % Adicionado 'anonymous' para a fase de revisão
\settopmatter{printacmref=false}
\renewcommand\footnotetextcopyrightpermission[1]{}
\setcopyright{none}

\acmConference[ISE 2026]{Workshop on Intelligent Software Engineering}{September 11, 2026}{São Paulo, SP, Brazil}

\AtBeginDocument{
    \providecommand\BibTeX{{Bib\TeX}}
}

\usepackage{graphicx}
\usepackage[brazilian]{babel}
\usepackage{pgfplots}
\pgfplotsset{compat=1.18}

\renewcommand\abstractname{Resumo}
\renewcommand\keywordsname{Palavras-chave}
\renewcommand\refname{Referências}

\renewcommand{\figurename}{Figura}
\renewcommand{\tablename}{Tabela}

\begin{document}

\title{Avaliação de LLMs Multimodais na Modalidade Programação da Olimpíada Brasileira de Informática: Uma Arquitetura Experimental}

% ------------------------------------------------------------
% NOTA: Mantenha os autores COMENTADOS durante a submissão!
% Descomente apenas quando o artigo for aprovado (Camera-ready)
% ------------------------------------------------------------
% \author{Victor Hugo Silva Ângelo}
% \affiliation{
%   \institution{Instituto da Computação - Universidade Federal de Alagoas}
%   \city{Maceió}
%   \country{Brasil}
% }
% \email{vhsa@ufal.ic.br}

\renewcommand{\shortauthors}{Ângelo}
\renewcommand{\shorttitle}{Resolução de Problemas da OBI com LLMs Multimodais}

\begin{abstract}
Este artigo tem como objetivo desenvolver e avaliar uma arquitetura para medir a capacidade de Grandes Modelos de Linguagem Multimodais (MLLMs) na resolução de desafios algorítmicos complexos na Modalidade Programação da Olimpíada Brasileira de Informática (OBI). Para isso, foi construída uma infraestrutura automatizada e experimental composta por um agente gerador de código baseado em LLM e um avaliador (Judge) para validação dos casos de teste. Dessa forma, um orquestrador gerencia o fluxo entre o banco de problemas, a estruturação de templates de prompts e os serviços externos (LLMs). Para simplificação dos resultados, adotou-se a seguinte configuração: o código gerado em Python e a técnica de prompt zero-shot, comparando-se enunciados contendo imagens no formato Base64 com os enunciados puramente textuais. O estudo traz um banco de questões ainda não explorado por MLLMs, além de traçar um paralelo pertinente entre a programação competitiva e a Engenharia de Software, no qual a interpretação de enunciados assemelha-se à modelagem de requisitos de sistemas e a validação de corner cases reflete a prática de testes unitários. Os resultados demonstram que o impacto da inclusão de imagens varia significativamente entre as ferramentas; embora o contexto visual possa adicionar ruído ou complexidade em alguns cenários, a análise comparativa revelou uma melhora no desempenho em problemas avançados de grafos e programação dinâmica em 3 dos 5 modelos testados, evidenciando que o suporte multimodal atua como um fator relevante para a resolução correta de desafios de alta complexidade algorítmica.
\end{abstract}

% Não esqueça de colocar suas palavras-chave reais aqui depois
\keywords{Modelos de Linguagem Multimodais, Geração Automática de Código, Engenharia de Software Experimental, Teste de Software, Programação Competitiva}

\frenchspacing
\maketitle

\section{Introdução}

A interpretação precisa de requisitos visuais é central na Engenharia de Software~\cite{Mendoza2024Comparative} e na programação competitiva. Nos problemas da Olimpíada Brasileira de Informática (OBI), ilustrações são frequentemente a chave para a compreensão de restrições algorítmicas, sem as quais a modelagem da solução torna-se incorreta.

Impulsionados pelas leis de escala~\cite{kaplan2020scaling}, os Grandes Modelos de Linguagem (LLMs) ganharam notoriedade na geração de código~\cite{Chen:2021:ELL} e na resolução de programação competitiva~\cite{Zheng2025LiveCodeBenchPH} e exames educacionais~\cite{Viegas2025AssessingTC, Mendona2024EvaluatingCV, nunes2023evaluating}. Com o avanço do \textit{Visual Instruction Tuning}~\cite{Liu2023VisualIT}, os Modelos de Linguagem Multimodais (MLLMs) ampliaram essa capacidade cognitiva para contextos com imagens~\cite{Fu2023MMEAC}, incluindo problemas matemáticos~\cite{Lu2023MathVistaEM}, avaliações universitárias~\cite{Almeida2023BLUEXAB, Yue2023MMMUAM} e, no ciclo de software, engenharia de requisitos~\cite{Mendoza2024Comparative}, testes~\cite{Silva2025LLMs} e correção de bugs visuais~\cite{Yang2024SWEbenchMD}.

Apesar desses avanços, o impacto do contexto visual no raciocínio algorítmico de nível educacional inicial permanece pouco explorado. A OBI é uma das principais portas de entrada ao pensamento computacional no Brasil, e seus enunciados dependem fortemente de figuras. Compreender como os MLLMs lidam com essas representações é relevante para a Inteligência Artificial e para a Engenharia de Software Experimental.

Neste trabalho, avaliamos uma arquitetura experimental automatizada para medir a capacidade de MLLMs na resolução de problemas da OBI~\cite{obi_passadas}. A infraestrutura integra um agente gerador e um juiz automático (\textit{Judge}) para validação com casos de teste oficiais. Para isolar o efeito visual, conduzimos experimentos \textit{zero-shot}~\cite{wei2021finetuned} em 181 questões, comparando enunciados textuais puros com versões contendo imagens (Base64) de cadernos da OBI entre 1999 e 2025.


\section{Trabalhos Relacionados}
Estudos recentes demonstraram a capacidade de Grandes Modelos de Linguagem na geração de código \cite{Chen:2021:ELL, Zheng2025LiveCodeBenchPH} e na resolução de exames educacionais \cite{Viegas2025AssessingTC, Mendona2024EvaluatingCV}. Com o avanço das técnicas de \textit{Visual Instruction Tuning} \cite{Liu2023VisualIT}, trabalhos têm avaliado Modelos de Linguagem Multimodais em raciocínio matemático com imagens \cite{Lu2023MathVistaEM} e tarefas complexas de Engenharia de Software \cite{Mendoza2024Comparative, Yang2024SWEbenchMD}. Destaca-se também a ideia de adicionar elementos visuais em enunciados de avaliações universitárias \cite{Almeida2023BLUEXAB}. No entanto, apesar desses avanços, o impacto e a eficácia de MLLMs no raciocínio algorítmico de nível educacional inicial, como os desafios da Olimpíada Brasileira de Informática (OBI), permanecem pouco explorados, motivando o presente estudo a avaliar a combinação de enunciados textuais e representações visuais.

\section{Metodologia}

Esta seção descreve o ambiente automatizado desenvolvido para avaliar a capacidade das MLLMs na resolução de problemas da modalidade Programação da OBI. O objetivo da metodologia é garantir um processo sistemático, reproduzível e comparável para análise do desempenho dos modelos.
Para reduzir variações no processo experimental, foi utilizado um formato padronizado de prompt, contendo:

\begin{itemize}
    \item O enunciado do problema com tags das imagens;
    \item A especificação de entrada e saída;
    \item Exemplos do enunciado;
    \item Instruções para gerar uma solução completa em código.
\end{itemize}

\subsection{Visão Geral da Arquitetura}

O sistema foi projetado para ser simples e automatizado, de forma que, ao escolher entre as linguagens de programação (C++ ou Python) e técnicas de prompt (Zero-Shot \cite{wei2021finetuned} ou Few-Shot \cite{brown2020language}), o sistema deve produzir um código que será averiguado por meio do gabarito. Assim, haverá um agente codificador (LLM Multimodal) e um avaliador (\textit{Judge}) para a dinamização do \textit{benchmarking} dos modelos utilizados para resolver os problemas da OBI de 1999 a 2025, disponíveis no site~\cite{obi_passadas}, conforme ilustrado na Figura~\ref{fig:arquitetura}.

\begin{figure}[ht]
  \centering
  \includegraphics[width=\linewidth]{samples/arquitetura-experimental-preto.png}
  \caption{Arquitetura Experimental para Resolução de Problemas de Programação da OBI.}
  \Description{Arquitetura Experimental}
  \label{fig:arquitetura}
\end{figure}

\subsubsection{Banco de Problemas}
\label{sec:Banco de Problemas}
\paragraph{Para a realização do experimento,}
foi construído um banco de problemas (\textit{dataset}) a partir do site oficial da Olimpíada Brasileira de Informática (OBI), utilizando \textit{script} em Python e convertendo os cadernos de questões na modalidade de programação em arquivos JSON por meio do modelo Gemini 3.1 Pro. As questões extraídas, as provas e os gabaritos estão disponíveis no seguinte repositório: \url{https://anonymous.4open.science/r/problems-of-the-obi-F517/README.md}.

\subsubsection{Camada Service}
\label{sec:Camada de Service}
\paragraph{Para o funcionamento do sistema,} foi necessário criar dois \textit{services} que se comunicam com o ambiente externo da arquitetura. O \textit{service} do LLM tem como objetivo realizar requisições à API das LLMs, calculando o custo total do prompt, além do total de tokens gastos e o código gerado pelas LLMs. Entretanto, o \textit{service} do \textit{Judge} tem como objetivo avaliar o código por meio de \textit{subthreads}, executando os casos de teste diretamente no ambiente de processamento. Com isso, o ambiente deve conter o \texttt{g++} e algum interpretador Python para garantir a automatização do processo. Além disso, o \textit{Judge} informa o status final da questão, a quantidade de acertos e de outros tipos de erros, e o tempo máximo de execução entre todos os casos de teste.

\subsubsection{Orquestrador e configuração}
\label{sec:Orquestrador e configuração}
\paragraph{Para garantir a execução contínua,} criou-se um orquestrador para executar no momento em que o usuário escolher uma configuração para execução da automatização. Esse orquestrador tem como objetivo executar as seguintes tarefas: criação e edição de arquivos; carregamento dos casos de teste da questão; validação do código; carregamento das imagens; e iteração com as questões da OBI para cada LLM.
Além disso, é importante a configuração de variáveis de ambiente para a conexão com as LLMs. Para garantir o funcionamento, é necessário informar: \texttt{NOME\_\_API\_KEY} (chave da API), \texttt{NOME\_\_BASE\_URL} (URL de comunicação), \texttt{NOME\_\_MODEL\_NAME} (nome do modelo) e os valores de custo por token (\texttt{NOME\_\_INPUT\_PRICE} e \texttt{NOME\_\_OUTPUT\_PRICE}).

\subsection{Modelos Avaliados}
Os experimentos utilizaram cinco modelos proprietários capazes de gerar código a partir de descrições textuais: Claude Sonnet 4.6 (Anthropic), Gemini 3.1 Pro (Google DeepMind), GPT 5.4 (OpenAI), Grok 4.20 (X) e Mistral Small 4.

\subsection{Ambiente de Avaliação}
O processo de avaliação das soluções geradas pelos modelos tem as seguintes etapas: seleção de um problema do banco de dados; envio do enunciado ao modelo de linguagem; recebimento do código gerado pelo modelo; compilação e execução do código; avaliação automática utilizando casos de teste.

A verificação das soluções é realizada comparando a saída produzida com o gabarito dos casos de teste. Cada solução recebe um veredito comum em juízes automáticos: \textit{Accepted} (AC, solução correta), \textit{Wrong Answer} (WA, saída incorreta), \textit{Compilation Error} (CE, erro de compilação), \textit{Runtime Error} (RE, erro de execução) ou \textit{Time Limit Exceeded} (TLE, excedeu o limite de tempo).

O limite de tempo para execução foi padronizado em 3 segundos para todas as instâncias e modelos deste estudo, pois essa informação não estava disponível nos cadernos de problemas da OBI.

\subsection{Configuração do Experimento}

No intuito de avaliar a capacidade das LLMs multimodais, foi necessário separar 181 questões que possuem imagens e casos de teste na base de problemas, sendo que essas questões compreendem 67 fáceis, 81 médias e 33 difíceis, segundo o Gemini 3.1 Pro.

Na configuração utilizada, o código deve ser produzido na linguagem Python com a técnica de \textit{prompt zero-shot}~\cite{wei2021finetuned}. Dessa forma, a imagem é inserida no prompt no formato Base64, com referências ao longo do enunciado, como ilustrado na Figura~\ref{fig:enunciado-vertical}.

Com isso, seguindo o pipeline ilustrado na Figura~\ref{fig:pipeline-questão}, para cada questão, o enunciado (com imagens em Base64, conforme Figura~\ref{fig:enunciado-vertical}) é inserido no template do prompt. O agente gera um arquivo em Python, e o \textit{Judge} valida os casos de teste salvando os resultados em um arquivo CSV.

\begin{figure*}[ht]
  \centering
  \begin{minipage}{0.48\textwidth}
    \centering
    \includegraphics[width=\linewidth]{samples/enunciado-prompt-horizontal.png}
    \caption{Formatação do Enunciado com Imagens em Base64.}
    \label{fig:enunciado-vertical}
  \end{minipage}\hfill
  \begin{minipage}{0.48\textwidth}
    \centering
    \includegraphics[width=\linewidth]{samples/pipeline.png}
    \caption{Pipeline de Execução para cada Questão.}
    \label{fig:pipeline-questão}
  \end{minipage}
\end{figure*}

Por fim, o ambiente deve ter acesso à internet para requisições, juntamente com um interpretador Python. O ambiente de \textit{hardware} é constituído por um processador Intel Core i7-1360P de 13ª geração (2,20 GHz), 16 GB de memória RAM (4800 MT/s) e acelerador gráfico integrado Intel Iris Xe, operando sob o sistema Windows 11 Home (64 bits).

\section{Resultados}
Esta seção tem como objetivo apresentar os resultados do experimento executado. Nas Subseções~\ref{subsec:resultados_sem_imagens} e \ref{subsec:resultados_com_imagens} serão mostradas as estatísticas gerais dos LLMs multimodais. Já a Subseção~\ref{subsec:desempenho_em_topicos} detalha os resultados para os casos de teste em tópicos avançados, como grafos e programação dinâmica.

\subsection{Resultados dos Enunciados Sem Imagens}
\label{subsec:resultados_sem_imagens}

O desempenho em relação à taxa de acerto para questões de nível fácil (acima de 77\%) e de nível médio (acima de 67\%) é considerado alto, como visto na Tabela~\ref{tab:desempenho_modelos_sem_imagens}, destacando-se os modelos Claude Sonnet 4.6, GPT 5.4 e Gemini 3.1 Pro. Entretanto, apenas o modelo Gemini 3.1 Pro obteve um desempenho satisfatório para questões consideradas difíceis. Vale ressaltar que este modelo também exigiu o maior tempo de processamento para a geração da resolução e o maior custo médio. Com isso, uma característica observada durante o experimento é que os modelos ainda têm dificuldades cognitivas para a resolução de problemas complexos. Um simples \textit{prompt} em configuração \textit{zero-shot} pedindo para resolver a questão não foi suficiente para passar no gabarito, pois uma característica intrínseca da programação competitiva é a necessidade de prever casos extremos (\textit{corner cases}) de entrada e saída.

\begin{table*}[htbp]
  \caption{Comparativo de Desempenho dos Modelos Multimodais sem Imagens}
  \label{tab:desempenho_modelos_sem_imagens}
  \begin{tabular}{l c c c c c c}
    \toprule
    \textbf{Modelo} & \textbf{Fácil AC (\%)} & \textbf{Médio AC (\%)} & \textbf{Difícil AC (\%)} & \textbf{Tokens I/O (Média)} & \textbf{Custo Médio (US\$)} & \textbf{Tempo Médio (s)} \\
    \midrule
    Claude Sonnet 4.6 & 77.61 & 74.07 & 36.36 & 1698 & 0.0121 & 9 \\
    GPT 5.4           & 80.60 & 67.90 & 27.27 & 1123 & 0.0059 & 4 \\
    Gemini 3.1 Pro    & 97.01 & 87.65 & 75.76 & 5128 & 0.0520 & 38 \\
    Grok 4.20         & 61.19 & 53.09 & 18.18 & 1958 & 0.0078 & 9 \\
    Mistral Small 4   & 58.21 & 48.15 & 9.09  & 1212 & 0.0016 & 3 \\
    \bottomrule
  \end{tabular}
\end{table*}

\subsection{Resultados dos Enunciados Com Imagens}
\label{subsec:resultados_com_imagens}

Ao aumentar o nível de complexidade do \textit{prompt} com a adição de imagens, apenas o modelo Gemini 3.1 Pro manteve os mesmos percentuais da Tabela~\ref{tab:desempenho_modelos_sem_imagens}. No entanto, ao analisar os demais modelos multimodais, observou-se um aumento na taxa de acerto para questões de nível fácil. Entretanto, para questões de nível médio, os modelos Gemini 3.1 Pro e Grok 4.20 mantiveram o mesmo desempenho, o GPT 5.4 obteve leve aumento, e os demais apresentaram queda na taxa de acerto. Além disso, para questões de nível difícil, apenas os modelos Gemini 3.1 Pro, GPT 5.4 e Mistral Small 4 conseguiram acertar mais questões utilizando o contexto visual. Isso evidencia que o aumento da complexidade do \textit{prompt} não é garantia de resolução para questões de níveis médio e difícil, como demonstrado na Tabela~\ref{tab:desempenho_modelos_com_imagens}, embora possa servir como um guia para a modelagem da solução.

\begin{table*}
  \caption{Comparativo de Desempenho dos Modelos Multimodais com Imagens}
  \label{tab:desempenho_modelos_com_imagens}
  \begin{tabular}{l c c c c c c}
    \toprule
    \textbf{Modelo} & \textbf{Fácil AC (\%)} & \textbf{Médio AC (\%)} & \textbf{Difícil AC (\%)} & \textbf{Tokens I/O (Média)} & \textbf{Custo Médio (US\$)} & \textbf{Tempo Médio (s)} \\
    \midrule
    Claude Sonnet 4.6 & 83.58 & 71.60 & 30.30 & 2049 & 0.0148 & 12 \\
    GPT 5.4           & 82.09 & 69.14 & 30.30 & 1329 & 0.0065 & 4 \\
    Gemini 3.1 Pro    & 97.01 & 87.65 & 78.79 & 6539 & 0.0567 & 39 \\
    Grok 4.20         & 73.13 & 53.09 & 15.15 & 1792 & 0.0061 & 7 \\
    Mistral Small 4   & 61.19 & 41.98 & 12.12 & 1456 & 0.0018 & 3 \\
    \bottomrule
  \end{tabular}
\end{table*}

\subsection{Desempenho em Tópicos Avançados}
\label{subsec:desempenho_em_topicos}

Cabe ressaltar que, para o \textit{Judge} aceitar uma submissão, o código deve passar em todos os casos de teste da questão. Ao analisar especificamente as questões que envolvem tópicos de grafos e programação dinâmica (64 questões e 2.491 casos de teste no total), é perceptível na Figura~\ref{fig:acuracia_modelos_total} que a utilização de imagens no enunciado proporcionou uma melhora no desempenho de acerto individual dos casos de teste. As exceções foram os modelos Claude Sonnet 4.6 e Mistral Small 4, que apresentaram um desempenho inferior ao processar o contexto visual.

\begin{figure}[H]
  \centering
  \begin{tikzpicture}
    \begin{axis}[
        ybar=0pt,
        width=\columnwidth,
        height=7cm,
        bar width=15pt,
        enlarge x limits=0.15,
        ymin=0, ymax=100,
        ytick={0, 20, 40, 60, 80, 100},
        yticklabel={\pgfmathprintnumber\tick\%},
        ylabel={\% de acertos nos casos de teste},
        ymajorgrids=true,
        grid style={line width=.1pt, draw=gray!20},
        axis x line*=bottom,
        axis y line*=left,
        symbolic x coords={Gemini 3.1 Pro, Claude Sonnet 4.6, GPT-5.4, Grok 4.20, Mistral Small 4},
        xtick=data,
        x tick label style={rotate=35, anchor=north east, font=\footnotesize},
        nodes near coords,
        nodes near coords style={font=\tiny, /pgf/number format/fixed, /pgf/number format/precision=1},
        legend style={
            at={(1,1)},
            anchor=north east,
            legend columns=-1,
            draw=none,
            font=\footnotesize
        },
    ]

    % Dados: Sem imagem (Vermelho)
    \addplot[fill=red!70!white, draw=none] coordinates {
        (Gemini 3.1 Pro, 76.3) 
        (Claude Sonnet 4.6, 53.5) 
        (GPT-5.4, 45.1) 
        (Grok 4.20, 37.6) 
        (Mistral Small 4, 29.3)
    };

    % Dados: Com imagem (Azul)
    \addplot[fill=blue!70!white, draw=none] coordinates {
        (Gemini 3.1 Pro, 81.0) 
        (Claude Sonnet 4.6, 51.1) 
        (GPT-5.4, 48.3) 
        (Grok 4.20, 38.0) 
        (Mistral Small 4, 23.5)
    };

    \legend{Sem imagem, Com imagem}
    \end{axis}
  \end{tikzpicture}
  \caption{Desempenho comparativo de acertos totais por modelo com e sem imagem.}
  \label{fig:acuracia_modelos_total}
  \Description{Gráfico de barras verticais mostrando que o Gemini 3.1 Pro lidera com 81.0 por cento no cenário com imagem e 76.3 por cento sem imagem, enquanto o Mistral Small 4 apresenta o menor desempenho, especialmente no cenário com imagem (23.5 por cento).}
\end{figure}

\section{Conclusão e Trabalhos Futuros}
\label{sec:conclusao}

Este trabalho apresentou uma arquitetura automatizada e um banco inédito de 493 questões para \textit{benchmarking} de MLLMs na OBI. Os experimentos com 181 problemas multimodais (\textit{zero-shot}, Python) indicam que a inclusão de imagens auxilia na compreensão semântica --- beneficiando questões de nível fácil e tópicos avançados na maioria dos modelos ---, mas não é suficiente para superar limitações em raciocínios algorítmicos complexos. A tensão entre o formato gerado pelos modelos probabilísticos e o determinismo do \textit{Judge} permanece como um desafio relevante na geração autônoma de código.

Como trabalhos futuros, sugerimos a replicação em \textit{datasets} diversificados para validar a capacidade de generalização e a experimentação com outras linguagens (como C++). A adoção de técnicas avançadas de \textit{prompting} (ex: \textit{Few-Shot} e \textit{Chain-of-Thought}) aliada a fluxos iterativos de \textit{feedback} com o \textit{Judge} são caminhos promissores para melhorar as taxas de acerto e aproximar a avaliação ao fluxo real de resolução de problemas de programação e Engenharia de Software.


\section*{Disponibilidade de Artefatos}

A arquitetura experimental, juntamente com os códigos e dados utilizados nas avaliações, encontra-se disponível para fins de revisão em: \url{https://anonymous.4open.science/r/obi-benchmarking-1472/README.md}

\bibliographystyle{ACM-Reference-Format}
\bibliography{sample-base}

\end{document}